% Conclusion -- handbook section 2.3 requires this coverage
% What the comparison establishes, the limits, and what comes next.

\section{Conclusion}

%----------------------------------------------------------
\begin{frame}{Answering the Research Question}
\footnotesize
\textbf{How much of the gap between a from-scratch model and a pretrained one comes from
pretrained features rather than architecture, loss or capacity?}

\begin{block}{A sharp answer}
Fully fine-tuned pretrained features are worth
\begin{itemize}
    \item $\mathbf{0.1301}$ validation mIoU over the from-scratch U-Net,
    \item $\mathbf{0.1175}$ over the \emph{same weights} with the backbone frozen,
    \item $\mathbf{0.1178}$ over the \emph{same model} trained with a class-balanced objective,
\end{itemize}
and those three alternatives are \textbf{mutually indistinguishable}, within 0.0126 of one
another. One factor accounts for essentially the whole spread.
\end{block}

\begin{alertblock}{An unexpected qualification}
The advantage is \textbf{not uniform}. The from-scratch U-Net is better, with separated intervals,
on \textbf{pole, sky and vegetation}, because a full-resolution decoder localises thin structures
that a stride-8 backbone cannot. Five of six error groups are boundary-dominated.

\textbf{Pretraining supplies recognition rather than resolution.}
\end{alertblock}

\vspace{0.1cm}
Carried forward: \textbf{DeepLabV3-MobileNetV3, fully fine-tuned}, at
$\mathbf{0.6180 \pm 0.0049}$ test mIoU and $0.7441 \pm 0.0048$ test Dice, 140\,s/epoch,
15.2\,ms/image.
\end{frame}

%----------------------------------------------------------
\begin{frame}{Limitations}
\small
\begin{itemize}
    \item \textbf{The U-Net is undertrained.} Validation mIoU still rising at epoch 40 on all
          three seeds, so $0.4836$ is a floor and the margin over it an \emph{upper bound}.
    \item \textbf{The model comparison confounds two factors.} Its arms differ in architecture \emph{and} in
          pretraining; no run here separates them. A DeepLabV3 trained from scratch would.
    \item \textbf{The backbone extension is one seed} and supports no claim of superiority.
    \item \textbf{Three seeds estimate a standard deviation from three observations.}
          ``Indistinguishable'' means \emph{not separated at this sample size}, not equal.
    \item \textbf{The data is narrow}: daylight, fair weather, European cities, one camera,
          and images were resized to a quarter of their pixels for compute, itself a ceiling on
          the thin classes that dominate the errors.
    \item \textbf{Reproducibility has a hardware limit.} The Apple GPU backend has no deterministic kernel for the
          interpolation DeepLabV3 uses, so rerunning a DeepLabV3 seed will not reproduce its
          number exactly. The reported spread includes that nondeterminism.
\end{itemize}
\end{frame}

%----------------------------------------------------------
\begin{frame}{What Comes Next}
\small
In the order the evidence supports them:

\begin{enumerate}
    \item \textbf{Resolution.} Train and evaluate at the native $1024 \times 2048$, or add a
          full-resolution decoder path (DeepLabV3+). This attacks the thin-structure and boundary
          groups, which are most of the error, and the $512 \times 1024$ resize was a choice made
          for compute rather than accuracy.
    \item \textbf{The clean version of the model comparison.} A DeepLabV3 trained from scratch separates pretraining
          from architecture in a single run.
    \item \textbf{The class-balancing question the ablation \emph{reopened}.} Dice accumulated over
          an epoch, or masking absent classes out of its mean, avoids the suppression mechanism
          measured here.
\end{enumerate}

\vspace{0.3cm}
\begin{block}{Two practices that changed the conclusions}
Reporting per-class numbers beside the mean, and fixing the decision rule before seeing results.
The per-class table located the one setting in which the \emph{weaker} model wins, and the rule
turned three apparently different outcomes into a single tie.
\end{block}
\end{frame}
